% =====================================================================
%  A4 プリント（profile: handout）用の追加プリアンブル。
%  pandoc の --standalone に --include-in-header で足される。
%  自分のものに差し替えるなら octavo.config.py に
%      'headers': {'latex': 'my-handout-header.tex'},
% =====================================================================

% ---- 日本語 -----------------------------------------------------------
% documentclass が ltjsarticle なら luatexja は既に読まれている。
% article（lang: 'en'）のときだけ luatexja-preset で和文を足す。
% texlive-lang-japanese が入っていない環境でも（日本語は化けるが）止まらない。
\makeatletter
\@ifpackageloaded{luatexja}
  {\IfFileExists{luatexja-fontspec.sty}{\usepackage{luatexja-fontspec}}{}}
  {\IfFileExists{luatexja-preset.sty}
     {\usepackage[haranoaji]{luatexja-preset}}
     {\PackageWarningNoLine{octavo}{luatexja が無い。日本語は組めない。
       sudo apt install texlive-lang-japanese (macOS: MacTeX)}}}
\makeatother

% ---- 版面 -------------------------------------------------------------
\usepackage[margin=22mm,top=20mm,bottom=20mm]{geometry}
\usepackage{enumitem}
\setlist{nosep,leftmargin=1.6em}
\usepackage{titlesec}
\titlespacing*{\section}{0pt}{1.2\baselineskip}{0.4\baselineskip}
\titlespacing*{\subsection}{0pt}{0.9\baselineskip}{0.3\baselineskip}

% 題扉を詰める（プリントは1行目から本文に入りたい）
\usepackage{titling}
\setlength{\droptitle}{-2.2\baselineskip}
\pretitle{\begin{center}\large\bfseries}
\posttitle{\end{center}\vspace{-0.6\baselineskip}}
\postauthor{\end{center}\vspace{-0.8\baselineskip}}
\postdate{\end{center}\vspace{-0.6\baselineskip}}

% ---- 図表 -------------------------------------------------------------
\usepackage{caption}
\captionsetup{font=small,labelfont=bf,skip=4pt}
\usepackage{adjustbox}
\let\oldtabular\tabular
\let\endoldtabular\endtabular
\renewenvironment{tabular}{\adjustbox{max width=\linewidth}\bgroup\oldtabular}
                          {\endoldtabular\egroup}
\newcommand{\inputtable}[1]{\IfFileExists{#1.tex}{\input{#1}}{\textbf{[表が無い: #1]}}}

% ---- 書き込み欄 -------------------------------------------------------
% 原稿で \fillin あるいは \fillin[8] と書くと下線が引ける（穴埋めプリント用）。
\newcommand{\fillin}[1][4]{\underline{\hspace{#1\zw}}}
% \zw が無い環境（article + luatexja-preset）向けの保険
\providecommand{\zw}{em}
